\section{Related Work}
\label{sec:related-work}

ViPragSent-Final sits between Vietnamese language modelling, affective-resource construction, and pragmatic-language evaluation. Prior work supplies model families, datasets, and task formulations, but none of the cited work by itself establishes the present six-label evaluation, the final ensemble's aggregate result, or a general advantage for a particular training strategy. We therefore separate related tasks from this paper's fixed evaluation and provenance-bounded contribution.

\subsection{Vietnamese Pretraining and Social-media Language}

PhoBERT provides a Vietnamese monolingual pretrained encoder, while XLM-R provides a multilingual representation-learning reference \citep{nguyen-tuan-nguyen-2020-phobert,conneau-etal-2020-unsupervised}. ViSoBERT extends this landscape with a Vietnamese social-media language-model context \citep{nguyen-etal-2023-visobert}. These studies motivate comparing monolingual, multilingual, and social-media-oriented encoders, but their pretraining objectives and benchmarks do not define the present pragmatic labels. Accordingly, ViPragSent evaluates these families under one fixed protocol rather than treating social-media pretraining as a guarantee of label-wise superiority.

\subsection{Vietnamese Affective Resources and Pragmatic Phenomena}

Vietnamese affective resources have principally targeted ordinary sentiment or emotion. UIT-VSFC is a student-feedback sentiment corpus, and UIT-VSMEC is a social-media emotion resource \citep{nguyen-etal-2018-uit-vsfc,ho-etal-2020-uit-vsmec}; ViGoEmotions further demonstrates fine-grained emotion classification on Vietnamese social-media comments \citep{tran-etal-2026-vigoemotions}. These resources establish useful task and domain precedents, but their label inventories do not jointly operationalize implicit sentiment, sarcasm, irony, figurative language, code-switching, and mocking. In this study they remain contextual or external-diagnostic resources, not sources of ViPragSent records.

Pragmatic interpretation is not reducible to ordinary polarity: verbal irony can depend on semantic incongruity and contextual strategies \citep{ghosh-etal-2019-verbal-irony}, while figurative-language evaluation requires its own linguistic and cultural coverage. The VIVID benchmark explicitly targets Vietnamese figurative-language gaps \citep{rem-lab-2026-vivid}. These lines of work motivate treating pragmatic labels as separate evaluation targets; they do not support a claim that the present corpus exhausts Vietnamese pragmatics. ViPragSent instead fixes a bounded six-label task and reports the provenance of its VIVID-derived candidate stratum separately.

\subsection{Social-media Variation and Code-switching}

Social-media-oriented models and emotion resources show that informal Vietnamese introduces domain-specific variation beyond conventional review or feedback data \citep{nguyen-etal-2023-visobert,ho-etal-2020-uit-vsmec,tran-etal-2026-vigoemotions}. Code-switching adds a further phenomenon whose evaluation cannot be inferred from ordinary sentiment scores alone. The present work therefore includes code-switching as one independently scored pragmatic label, rather than using a generic social-media claim as a substitute for evidence. This design leaves broader questions about cross-domain generalization and code-switching mechanisms open.

\subsection{Auxiliary Objectives, Ensembles, and Leakage-safe Selection}

Multi-task work has shown that related linguistic objectives can be jointly modelled with targeted sentiment classification \citep{moore-barnes-2021-multi}. That precedent motivates evaluating auxiliary objectives, but it does not imply that a fixed auxiliary set benefits a new corpus or backbone. ViPragSent therefore preserves its historical ablations as configuration-specific evidence and does not attribute the promoted result to any single auxiliary task.

The promoted system is not presented as a novel ensemble architecture. It is a frozen label-wise composition selected with five-fold development out-of-fold evidence and evaluated once after freeze on canonical test labels. This protocol distinguishes development selection from test evaluation; it is not a claim that no prior work uses specialization or ensembling. Its contribution is the documented, leakage-bounded application to the fixed ViPragSent task, including transparent aggregate and per-label comparisons.

\subsection{Instruction-tuned Baselines in Low-resource Evaluation}

Instruction-tuned and prompted models provide useful comparison families, but their behaviour in lower-resource settings is empirical rather than assumed. A controlled cross-lingual study compares prompting, translation, fine-tuning, and instruction tuning across target languages and downstream tasks \citep{toukmaji-flanigan-2025-adapting}; Sailor supplies an open South-East Asian language-model family \citep{dou-etal-2024-sailor}, while Vistral supplies the Vietnamese instruction-tuned model card used for the configured baseline \citep{nguyen-etal-2023-vistral}. These sources motivate reporting prompted and 7B instruction-tuned baselines with their recorded scope. They do not establish that an LLM family must win Vietnamese pragmatic detection, nor that a repository-specific QLoRA run is independently reproducible.

For AIVIVN, Nguyen et al. provide a proceedings description of a binary Vietnamese e-commerce-review setting \citep{nguyen-etal-2020-efficient}. We cite it only for that published setting: the repository's evaluated artifact is a separately hash-identified Kaggle mirror used for external polarity diagnostics, not an asserted canonical organizer release. Across these strands, the gap addressed here is a provenance-bounded comparative evaluation of six Vietnamese pragmatic phenomena, with a development-selected final system whose aggregate lead is explicitly separated from per-label leadership.
